In a reversible dynamics no information is destroyed, so the record of an interaction must be
somewhere afterwards. We ask where. Using the 256 second-order (Fredkin) extensions of the
elementary cellular automata, $x_{t+1}=f(x_t)\oplus x_{t-1}$, we first identify bounded, autonomous,
slower-than-light coherences by frozen, outcome-blind searches, and then collide them. For each
collision we compute an exact ledger over 64 impact offsets: how much of the information that
distinguishes the inputs leaves with the outgoing coherences, and how much stays behind in the
background. We then intervene: every resolved coherence is erased, and a fresh probe coherence is
launched through the remaining background. Across \nsthreePairs{} collision pairs in 25 rules, all
\nsthreeClean{} pairs whose background returns exactly to vacuum give zero information to the probe,
while in all \nsthreeShadow{} pairs that leave a residue (\nsthreeShadowRules{} rules) the probe's effect on
the future recovers \emph{all} of the information held only by the background
(\nonlyMin--\nonlyMax{} bits). In \nsthreeStatePos{} of those pairs the outgoing probe coherence itself carries
part of the record away, and in \nsthreeStateOne{} it carries all of it. Every rule that leaves such a
``shadow'' has an oscillating vacuum. All tests were specified and frozen before their outcomes
were computed. Invalid runs and falsified hypotheses are reported alongside the successes.
